\documentclass{report}
\usepackage{graphicx}
\usepackage{geometry}
\usepackage{caption}
\usepackage{subcaption}
\usepackage{amssymb}
\usepackage{amsmath}
\usepackage{enumitem}
\usepackage{physics}
\usepackage{hyperref}
\usepackage{cleveref}
\usepackage{url}
\usepackage{dirtytalk}
\usepackage{color}
\usepackage{calc}
\usepackage{float}
\usepackage{epsfig}
\usepackage{dsfont}
\usepackage{lipsum}
\usepackage{tikz-feynman}
\usepackage{multirow}
\usepackage{makecell}
\usepackage{MnSymbol}%\sumint
\usepackage{wrapfig}
\usepackage{textpos}
\usepackage{biblatex}
\usepackage[integrals]{wasysym}
\addbibresource{bibliography/bibliography.bib}
\usepackage{fancyhdr}
%custom page style for report class
\fancypagestyle{mystyle}{
\pagestyle{fancy}
\fancyhf{} % clear all headers      
\fancyhead[R]{\rmfamily \small \nouppercase \rightmark}
\fancyhead[L]{\rmfamily \small \nouppercase \leftmark}  
\fancyfoot[C]{\thepage}
\renewcommand{\headrulewidth}{0pt}
\renewcommand{\footrulewidth}{0pt}}
\newcommand{\alert}[1]{{\color{red} #1}}


% macros for fast scaling parentheses
\newcommand{\lp}{\left(}
\newcommand{\rp}{\right)}

\newcommand{\lsb}{\left[}
\newcommand{\rsb}{\right]}

\newcommand{\lcb}{\left\{}
\newcommand{\rcb}{\right\}}

\newcommand{\la}{\left\langle}
\newcommand{\ra}{\right\rangle}
\definecolor{red}{rgb}{0.631, 0.094, 0.094}
\begin{document}

%%% TITLE PAGE %%%

\thispagestyle{empty}
\begin{textblock*}{\textwidth}(-2cm,-2.5cm)
\includegraphics[height=0.1\paperheight]{logos/Bicocca_transparent.png}	
\end{textblock*}
\begin{textblock*}{\textwidth}(1.9cm,-2.6cm)
\begin{flushleft}
{\color{red}\Large{SCUOLA DI DOTTORATO}}\\[2mm]
{\color{red}\large{UNIVERSIT\`A DEGLI STUDI DI MILANO-BICOCCA}}
\end{flushleft}
\end{textblock*}
\vspace{2cm}
\centerline{\Large{Department of Physics \say{G. Occhialini}}}
\vspace{0.5cm}
\centerline{\Large{Ph. D. program in Physics and Astronomy, XVIII cycle}}
\vspace{2mm}
\centerline{\Large{Curriculum of Theoretical Physics}}
\vspace{2cm}
\centerline{\huge{ The non-static screening spectrum}}
\centerline{\huge{of high-temperature Quantum Chromo-Dynamics}}
\vspace{2cm}
\centerline{\LARGE{Pietro Rescigno}}
\vspace{2mm}
\centerline{\Large{Registration number: 826333}}
\vspace{2cm}
\begin{flushleft}
\LARGE{Tutor: Leonardo Giusti}\\[1.5mm]
\LARGE{Supervisor: Marco Cè}\\[3cm]
\LARGE{Coordinator: Stefano Ragazzi}
\end{flushleft}
\vfill

\centerline{\LARGE{Academic Year 2025/2026}}
%%%%%%%%%%%%%%%%%
%%% ABSTRACT (Techincal)%%%%
\begin{abstract}
   A one-page \say{technical} abstract that summarizes
   the research work, similar to the one in the 
   paper on baryonic screening masses.	
\end{abstract}
%%% PUBLICATIONS %%%
\pagenumbering{roman}

{\centering
{\Large\textbf{Scientific Output}}\\[5mm]
\centerline{Here I list the works I have published and 
my contributions to}
\centerline{international conferences, either
as a speaker or as an author.}
\vspace{3mm}
\centerline{Bibliometric information can be retrieved on my}
\centerline{\href{https://inspirehep.net/authors/2727845?ui-citation-summary=true}{iNSPIRE} or \href{https://scholar.google.com/citations?user=3bmCuiAAAAAJ&hl=it}{Google Scholar} pages.}}


% Name & headline

%\switchcolumn
\section*{List of publications}
$\bullet$ L. Giusti, T. Harris, D. Laudicina, M. Pepe, PR, \textit{Baryonic screening masses in QCD at high temperature}, Phys. Lett. B 855 (2024) 138799,
\href{https://doi.org/10.1016/j.physletb.2024.138799}{https://doi.org/10.1016/j.physletb.2024.138799}, \href{https://arxiv.org/abs/2405.04182}{arxiv:2405.04182}\\[5mm]
$\bullet$ L. Giusti, M. Laine, D. Laudicina, M. Pepe, PR, \textit{Baryonic thermal screening mass at NLO}, JHEP 06 (2024) 205, \href{https://doi.org/10.1007/JHEP06(2024)205}{https://doi.org/10.1007/JHEP06(2024)205}, \href{https://arxiv.org/abs/2405.03975}{arxiv:2405.03975}

\section*{Contributions to conferences}
As a \textbf{speaker:}
\begin{itemize}
\item PR, L. Giusti, T. Harris, M. Laine, D. Laudicina, M. Pepe, \textit{Nucleon screening 
masses in QCD at high temperature}, talk - Strong and
Electroweak Matter 2024, Goethe Universit\"at Frankfurt, DE - \href{https://indico.physik.uni-bielefeld.de/event/100/contributions/226/}{Indico page} \\
\item PR, L. Giusti, T. Harris, D. Laudicina, M. Pepe, \textit{Baryonic screening masses in 
high temperature QCD}, talk -
The 40$^{\rm th}$ International Symposium on Lattice Field Theory, Fermilab, USA - \href{https://pos.sissa.it/453/196/}{PoS(LATTICE2023)}
\end{itemize}
As an \textbf{author:}
\begin{itemize}
    \item L. Giusti, M. Bresciani, M. Dalla Brida, T. Harris, D. Laudicina, M. Pepe,
    PR, \textit{Non-perturbative thermal QCD at very high temperatures: computational strategy and hadronic screening masses}, talk - The 41$^{\rm st}$ International Symposium on Lattice Field Theory, University of Liverpool, UK - to appear in 
    \textit{Proceedings of Science} - \href{https://conference.ippp.dur.ac.uk/event/1265/contributions/7291/}{Indico page}
    \item L. Giusti, M. Bresciani, M. Dalla Brida, T. Harris, D. Laudicina, M. Pepe,
    PR, \textit{Non-perturbative thermal QCD at 
    very high temperatures}, 
    42$^{\rm nd}$ International Conference on High Energy Physics, Prague, CZ - to appear in \textit{Proceedings
    of Science} - \href{https://arxiv.org/abs/2410.12626}{arXiv:2410.12626}
\end{itemize}
%%%%%%%%
\tableofcontents
\pagestyle{mystyle}

\addcontentsline{toc}{chapter}{Introduction}
\chapter*{Introduction}
A somewhat concise ($\sim$ 3 pages) introduction about your research


\pagenumbering{arabic}

%\pagestyle{mystyle}
\chapter{Name of the chapter}

Text here


\chapter{Quantum Chromo-Dynamics on the lattice}
\section{Lattice QCD at $T=0$}
Definition of QCD on a finite lattice: Wilson gauge action, Wilson fermions,
doublers + Nielsen-Ninomiya, symmetries of the lattice/continuum actions, 
path integral, Monte Carlo and Metropolis algorithm,
correlation functions, spectrum, systematic effects (discretization, finite volume),
scale setting, renormalization, continuum limit\\
\section{Lattice QCD at $T\neq 0$}
Usual discussion


\chapter{Screening masses}
\section{Definition}
Definition of screening masses, screening correlators, relation with 0-temp.
mass.
\section{Analytic results}
Review of analytically known results (EFT, mesons, baryons)
\section{Lattice results}
Review of lattice determinations (mesons, baryons)
\subsection{JHEP '22}
Dedicated section on meson paper of '22: non-perturbative high-temperature 
strategy, comparison with PT, splitting\\
\section{Aim of the thesis}
Conclude chapter with the \alert{aim of the thesis:} compute baryonic screening
masses from the lattice at high T, and at NLO in EFT\\
Tackle exponential problem in non-static mesons, detail spin-splitting


\chapter{Original contributions}

Detail the original work I carried during Ph. D.
\section{Baryonic screening masses}
\subsection{NLO EFT}
Describe \cite{BaryonNlo}
\subsection{From the lattice}
Describe \cite{BaryonLattice} 
\section{Non-static mesonic screening masses}
Describe SNR problem, detail random sources implementation, determine screening
spectrum
\section{Hyperfine splitting in static mesons}
\subsection{At order $g^4$ in EFT perturbation theory}
\subsection{With high-precision data from lattice simulations}


\addcontentsline{toc}{chapter}{Conclusions and Outlook}
\chapter*{Conclusions and Outlook}


\appendix
\chapter{Notation and conventions}\label{app:A}


%%%% BIBLIOGRAPHY %%%%%
\addcontentsline{toc}{chapter}{Bibliography}
\printbibliography

\addcontentsline{toc}{chapter}{Acknowledgements}
\chapter*{Acknowledgements}

\end{document}

